% Generated draft: requires mathematical and source review.
\documentclass[11pt]{article}
\usepackage[margin=25mm]{geometry}
\usepackage{fontspec}
\setmainfont{Latin Modern Roman}
\usepackage{amsmath,amssymb,mathtools,unicode-math}
\setmathfont{Latin Modern Math}
\usepackage{enumitem,needspace,xurl,hyperref}
\setlength{\parindent}{0pt}
\setlength{\parskip}{5pt}
\setlength{\emergencystretch}{3em}
\newcommand{\N}{\mathbb{N}}
\newcommand{\Z}{\mathbb{Z}}
\newcommand{\Q}{\mathbb{Q}}
\newcommand{\R}{\mathbb{R}}
\newcommand{\C}{\mathbb{C}}
\newcommand{\F}{\mathbb{F}}
\newcommand{\A}{\mathbb{A}}
\newcommand{\PP}{\mathbb{P}}
\newcommand{\E}{\mathbb{E}}
\newcommand{\eps}{\varepsilon}
\newcommand{\dd}{\,\mathrm{d}}
\begin{document}
\section*{UnsolvedMath ID: 701531}
\textbf{Source problem code:} catalogue entry
\subsection{Definitions and mathematical statement}
A group \(G\) is \emph{minimax} if it has a finite subnormal series
\[
1=G_0\trianglelefteq G_1\trianglelefteq\cdots\trianglelefteq G_n=G
\]
such that each factor \(G_i/G_{i-1}\) satisfies either the minimum condition on subgroups (every descending chain of subgroups stabilizes) or the maximum condition on subgroups (every ascending chain of subgroups stabilizes). The choice of condition may differ between factors. A group is \emph{nilpotent} if its lower central series, defined by \(\gamma_1(G)=G\) and \(\gamma_{j+1}(G)=[\gamma_j(G),G]\), reaches the trivial subgroup after finitely many steps.

A finitely generated field of characteristic zero is a field \(K=\Q(a_1,\ldots,a_r)\) for some finite list of elements. For a nonzero \(K\)-vector space \(V\), a representation \(\rho:G\to\operatorname{GL}_K(V)\) is \emph{faithful} if \(\ker\rho=\{1\}\), and \emph{irreducible} if its only invariant \(K\)-linear subspaces are \(0\) and \(V\). Using the algebraic imprimitivity convention, it is \emph{primitive} if there is no decomposition
\[
V=\bigoplus_{i\in I}V_i,\qquad |I|>1,
\]
into nonzero \(K\)-linear subspaces that \(\rho(G)\) permutes transitively.

\textbf{Question (A. V. Tushev; Kourovka Notebook 18.103).} Is it true that for every nilpotent minimax group \(G\), every finitely generated field \(K\) of characteristic zero, and every nonzero \(K\)-vector space \(V\),
\[
\left(\exists\,\rho:G\longrightarrow\operatorname{GL}_K(V)\text{ faithful, irreducible, and primitive}\right)
\quad\Longrightarrow\quad
\left(\exists\,S\subseteq G\text{ finite with }G=\langle S\rangle\right)?
\]
\textbf{Unresolved qualification.} The supplied source does not specify the dimension of the representation space or define primitive representations. The displayed formulation allows arbitrary dimension and uses the algebraic direct-sum convention above; these conventions require confirmation against the original source.
\subsection{Short English statement}
Must a nilpotent minimax group admitting a faithful irreducible primitive representation over a finitely generated field of characteristic zero be finitely generated? The source leaves the representation dimension and the precise meaning of primitive unspecified; the formulation above makes its conventions explicit.
\subsection{Sources}
\begin{itemize}
\item Supplied local UnsolvedMath registry record, permanent ID \(701531\), exact source code \texttt{catalogue entry}, titled \emph{Kourovka Notebook Problem 18.103}. The registry labels the problem open; this supplied label has not been independently verified.
\item \emph{The Kourovka Notebook (Unsolved Problems in Group Theory)}, original problem 18.103, proposed by A. V. Tushev in 2014. The supplied attribution names the 18th issue (2014) and the 21st issue edition (2026), with identifier \texttt{arXiv:1401.0300v48}. The collection record names E. I. Khukhro and V. D. Mazurov as editors. These bibliographic details are supplied attribution; the referenced source has not been independently consulted.
\end{itemize}
\end{document}
